\documentclass[12pt]{article}
\usepackage[a3paper, margin=1in]{geometry}
\usepackage{amsmath, amssymb, amsthm, ulem, booktabs}

\begin{document}

\section*{Domácí úkol 5.1}

Vypracovali: Daniel \textit{``Localhost"} Ransdorf, Jan \textit{``kokeshh"} Kodeš \hfill Podpis: \rule{4cm}{0.4pt}


\bigskip

Před důkazem nejprve vyslovme a dokažme jednoduché lemma.

\textbf{Lemma $(*)$}: Necht' $B$ je matice nad T taková, že součet prvků v každém řádku je 0.
Matice $B$ má potom netriviální jádro.

\begin{quotation}
  \textbf{Důkaz}: Pro každý řádek matice $B$ platí, že sečtením všech prvků na daném řádků získáme 0.
    Tedy součet všech sloupců je roven nulovému vektoru.
    Z toho plyne, že homogenní soustava $B\vec{x} = \vec{0}$ má netriviální řešení $\vec{x}_1 = (1,1,1,...,1)^T$.
    Tedy jádro matice $B$ je netriviální, protože obsahuje nenulový vektor $(1,1,1,...,1)^T$. \quad $\square$
\end{quotation}

\textbf{Důkaz zadaného tvrzení}:
Mějme čtvercovou matici $A$ a skalár $a$ ze zadání. Dle tvrzení ze skript je skalár $a$ vlastním
číslem matice $A$ právě tehdy, když $Ker\left(A-aI_n\right) \neq \left\{\vec{0}\right\}$.
Dále využijme ekvivalencí z kapitoly o determinantech:

\begin{align*}
  &Ker\left(A-aI_n\right) \neq \left\{\vec{0}\right\} \\
  \Longleftrightarrow \quad
  &det\left(A-aI_n\right) = 0 \\
  \Longleftrightarrow \quad
  &det\left(A-aI_n\right) = det\left(A-aI_n\right)^T
  = det\left(A^T-aI_n^T\right) = det\left(A^T-aI_n\right) = 0 \\
  \Longleftrightarrow \quad
  &Ker\left(A^T-aI_n\right) \neq \left\{\vec{0}\right\}
\end{align*}

Spojením dostáváme, že $a$ je vlastním číslem $A$ právě tehdy, když je jádro $(A^T-aI_n)$ netriviální. ← označme $(**)$

Součet prvků v každém sloupci $A$ je roven $a$. Tedy součet prvků v každém řádků $A^T$ je také roven $a$.
Nahlédněme, že odečtením $aI_n$ zmenšujeme součet na každém řádku o hodnotu $a$.
Tedy součet prvků v každém řádku $(A^T-aI_n)$ je roven $a - a = 0$.
Podle lemmatu $(*)$ máme, že jádro $(A^T-aI_n)$ je netriviální.

Z $(**)$ finálně plyne, že $a$ je vlastním číslem matice $A$. \quad $\square$
  
\end{document}
